\documentclass[11pt,a4paper]{article}
\usepackage[T1]{fontenc}
\usepackage{lmodern,amsmath,amssymb,amsthm,mathtools,booktabs,microtype,xcolor}
\usepackage[margin=26mm]{geometry}
\usepackage[colorlinks=true,linkcolor=blue!45!black,citecolor=blue!45!black,urlcolor=blue!45!black]{hyperref}
\newtheorem{theorem}{Theorem}[section]
\newtheorem{lemma}[theorem]{Lemma}
\newtheorem{proposition}[theorem]{Proposition}
\theoremstyle{remark}\newtheorem{remark}[theorem]{Remark}
\newcommand{\Q}{\mathbb Q}\newcommand{\Z}{\mathbb Z}\newcommand{\HH}{\mathcal H}
\DeclareMathOperator{\Tr}{Tr}
\setlength{\emergencystretch}{2em}
\title{A family of quartics in 297 variables}
\author{}\date{9 November 2026}
\begin{document}\maketitle
\begin{abstract}
We specify a family of integer quartics in 297 total variables by 271
quadratic equations and complete coefficient tables. Under the arithmetic
premises of the elliptic source construction, every sufficiently large
member defines the integers in the rationals. Two height detectors use
the rational trace of a sum of elliptic abscissas. The sixteen local tests
are represented by polynomial expressions with distinct valuation classes,
using fourth and eighth powers already computed for the formal logarithm.
The coefficient family is explicit, but the sufficient parameter threshold
is ineffective because the aggregate-trace height theorem has an
ineffective additive constant. No particular defining member of this
family is selected. A fully specified 342-variable quartic is a separate
result.
\end{abstract}

\section{The exact family and its scope}
All variable counts include the input $X$. Put
\[
                     A_N=80989^{2N+1}\qquad(N\ge1).
\]
The tables \nolinkurl{quadratic-family297.json} and
\nolinkurl{polynomial-family297.json} specify 271 primitive integer
quadratic rows $Q_{N,j}$ in 297 coordinates and
\[
                           P_N=\sum_{j=1}^{271}Q_{N,j}^2.
\]
Neither $A_N$ nor $N$ is a variable of $P_N$.

\begin{theorem}\label{thm:main}
Under the arithmetic premises in Section~\ref{sec:premises}, there is
an integer $N_0\ge27464$ such that, for every $N\ge N_0$ and $x\in\Q$,
\[
 x\in\Z\quad\Longleftrightarrow\quad
 \exists\boldsymbol y\in\Q^{296},\ P_N(x,\boldsymbol y)=0.
\]
For every $N\ge1$, $P_N$ has degree exactly four, content one, and
coefficient height
\[
                           H(P_N)<2^{1520}80989^{4N+2}.
\]
For every $N\ge1$ it has 4268 nonzero coordinate monomials.
Its exact height is $80989^{4N+2}+1$ for $N\ge23$.
The proof does not determine a numerical value of $N_0$.
\end{theorem}

An exact family specification and an effective choice of a defining
member are different assertions. The theorem supplies the first together
with eventual definition of $\Z$; its height argument does not supply
the second. The explicit 342-variable theorem in \cite{Explicit} uses
three height detectors and an effective sum-height inequality.

\section{Arithmetic premises}\label{sec:premises}
Let $K=\Q(\rho)$, $\rho^2=2$, $d=3-\rho$, and
\[
 E:y^2=x^3-d^2x,\qquad
 P=2(1+\rho,2\rho)
   =\left(\frac{57-28\rho}{8},\frac{-660+417\rho}{32}\right).
\]
The arithmetic input contract is the one precisely stated in
\cite[Section~2]{Source} and restated in \cite[Section~2]{Explicit}:
\begin{enumerate}
\item $2E(K)=\Z P$ and $h(x([n]P))\ge n^2-6$ for every nonzero
integer $n$.
\item The five-contact height bound $h(t)\le\HH M_5(t)^{19}$, with
\[
 \HH=2^{110}210^{1800}\left(\prod_{p\le76176}p\right)^{19}+128
       <2^{2093448}.
\]
Here $M_5(u/v)$ is the product of primes greater than $76176$ with
positive odd valuation on at least one of
$v,u,u-v,u-2401v,32u+343v$, in lowest terms with $v>0$, and the
excluded contact values are $0,1,2401,-343/32$.
\item The source pole obstruction outside $2,3,17$ and its simultaneous
realization for inputs in $(0,17)$ with square reduced denominator and
the specified local squareclasses and valuation bounds at $2,3,17$.
The normalized seven-coordinate equations are retained literally.
\item The simultaneous prime-pattern and conjugate reduction-order
conditions with arbitrary fixed odd avoidance modulus, finite congruences,
and real approximation, including the quantitative weighted selection
premise stated in the cited contract. In particular the slowly decreasing
separation parameter is available after all fixed data are chosen.
\end{enumerate}
No finite coefficient computation proves these inputs. The weighted selection and
separation used here follow from the cited analytic estimates as proved
in \cite{Selection}; no additional prime-distribution conjecture is
introduced by the local weight replacement. The
aggregate-trace height theorem used next is a proved consequence of
Von Buhren's family height theorem, with an ineffective additive constant,
as detailed in \cite{Aggregate}. It is not an additional conjecture.

\section{The trace height and the two local detectors}
The retained square charts give root points
\[
 Q_1=[n]P,\quad Q_2=[a]P,\quad Q_3=[t]P,\quad Q_4=[b]P,
 \qquad B=\max(|n|,|a|,|t|,|b|),
\]
with nonzero integer labels. Their interpretation for all nonzero
integer labels, without a parity restriction, is proved in
\cite[Section~2]{Baseline}. Define
\[
 S_0=\frac12\Tr_{K/\Q}\left(\sum_{i=1}^4x(Q_i)\right).
\]
There are fixed constants $M_1>0,M_2\ge1$ such that
\begin{equation}\label{eq:traceheight}
                           B^2\le M_1h(S_0)+M_2.
\end{equation}
The additive constant is not effective in the cited proof. Every root
abscissa is positive at both real embeddings, so $S_0>12$.

Let $\sigma$ denote conjugation on $K/\Q$ and let $f=1/x$, extended
by zero at the identity. The auxiliary comparison points satisfy
$U_0=H$, $V_0=H-P$, and $Q_i=U_i-V_i$ as group differences on $E$.
These are the exact addition relations of the retained point charts
in \cite{Baseline}.

Use the two nonzero rational square scales $1,S_0^2$. For each
scale $a_*$ put $t_*=a_*\lambda$, $\delta=t_*-1$ and
\[
 \begin{gathered}
 \Theta=\{0,-1/2,-1,1/2399,-32/407\},\\
 R_* =\frac{\prod_{\theta\in\Theta}(t_*-\theta)}{a_*^5\delta^6},
 \qquad q=R_*^{-1}\delta^{-4},\qquad q_m=q^{2m-1}.
 \end{gathered}
\]
These formulas are used on their nonzero domain; the exact system
consists of their cleared quadratic equations. The case $t_*=1$
has empty contact radical and requires no reciprocal.

Define $A_i=f(U_i)$, $B_i=\sigma f(V_i)$, and
\[
 \begin{aligned}
 C_i&=L_i-L'_i\quad(1\le i\le4),& I&=L_1-\beta,\\
 D_H&=f(H),&G&=L_3^2-L_2(L_2-2)L_3L_4+(L_2-2)^2L_4^2-1.
 \end{aligned}
\]
The alternate ratios $L'_i$ already include conjugation. The nine tests
$A_i,B_i$ have order one, $C_1,I$ have order ten, and $D_H,C_2,C_3,C_4,G$
have order forty. For an ordered list of tests, use the factor
\[
 \mathcal T(D_0,\ldots,D_{s-1})
 =1+\sum_{j=0}^{s-1}q_{m(D_j)}^{2^j}D_j^{2^{j+1}},
\]
where $m(D_j)$ is its stated order. Thus different required orders can
occur within one factor.

The detector uses the product of the factors for the following six lists:
\[
 \begin{gathered}
 (D_H,C_1,B_0),\quad(C_2,I,B_2),\quad(C_3,A_2,B_3),\\
 (C_4,A_3,B_4),\quad(G,A_4),\quad(A_1,B_1).
 \end{gathered}
\]
Each of the sixteen tests occurs exactly once. The common-place valuation
argument and the literal cleared trace equations are proved in
\cite[Proposition~2.3]{Shared}; each contact prime supplies every required test order
at one place. The rational logarithm approximants and the scalar point
exclusions are retained.

If a nonzero integer $J$ is divisible by the $k$th powers of both
contact radicals, the contact-height theorem bounds
$h(\lambda)$ and $h(S_0^2\lambda)$ by
$(\HH+\log3)|J|^{19/k}$. Since $\lambda\ne0$,
\[
 2h(S_0)=h(S_0^2)\le h(S_0^2\lambda)+h(\lambda).
\]
Combining this with \eqref{eq:traceheight} gives
\begin{equation}\label{eq:height}
 B^2\le C|J|^{19/k}+M_2,
 \qquad C=M_1(\HH+\log3).
\end{equation}


\subsection{A quadratic presentation of the local weights}
For each detector, write $F_*(t_*)=\prod_{\theta\in\Theta}(t_*-\theta)$
and put $q=R_*^{-1}\delta^{-4}$. The weights used here
are $q_m=q^{2m-1}$, so $q_{10}=q^{19}$ and $q_{40}=q^{79}$.
The quadratic rows are
\[
 \begin{gathered}
 q_2=q^2,\quad q_4=q_2^2,\quad c q_4=q,\quad q_8=q_4^2,\\
 c q_{19}=q_8^2,\quad q_{38}=q_{19}^2,\quad c q_{79}=q_{38}^2.
 \end{gathered}
\]
At every consumed contact prime, $\delta$ is a unit and $q\ne0$.
The rows then force $c=q^{-3}$ and each indicated power uniquely.
No inverse coordinate other than $c$ is needed. If $q=0$, the retained
relations $\delta^3q=e$ and $F_*(t_*)e=(a_*\delta)^5$ force
$\delta=0$; the contact radical is empty in that case.
This covers the potentially degenerate stratum without using a rational
formula outside its domain. The full local extraction proof from the
literal trace equations is \cite[Section 2]{Shared}.

At contacts these powers have exactly the required valuations:
$v(q_m)=(1-2m)v_p(R_*)$. The ordered factors therefore supply the same
orders $1,10,40$. At prepared finite primes their simultaneous smallness
is established directly in \cite[Section 3]{Shared}, by enlarging a fixed
denominator before prime selection. Thus the replacement changes neither
the formal logarithm comparisons nor their required depths.


\subsection{Primitive normalization of the scalar contact product}
Let $z_j$ denote a coordinate by its source identifier. The actual
contact product is the affine expression
\[
 \begin{aligned}
 r={}&-103171881284-2103220824z_{248}-21050600z_{487}\\
     &-7629700z_{488}-45771136z_{489}.
 \end{aligned}
\]
We use $r$ itself as the coordinate in place of $z_{489}$. This is an
invertible rational affine change: its coefficient of $z_{489}$ is
nonzero, and every other coordinate, including the input, is fixed.
The source coordinate occurs only in its own defining row and the
scalar normalizer row. After substitution and primitive normalization,
the latter is exactly
\[
                  R_s+80989^2rR_s-A_Nr=0.
\]
Every equation remains quadratic. Reverse substitution reproduces both
source rows exactly up to a nonzero rational scalar, so their rational
solution sets correspond bijectively. Taking sums of squares preserves
the required conjunction over $\Q$.

In the resulting quartic the coefficient of $r^2$ is $A_N^2+1$.
Every other coefficient is either a nonzero constant or a nonzero constant
multiple of $A_N$. Thus no coefficient vanishes for a positive parameter,
and the support counts hold for every $N\ge1$.
Let $C_0$ be the largest absolute constant coefficient. Its exact integer
value is recorded in the certificate and has 1520 binary digits. The
largest absolute coefficient multiplying $A_N$ is
\[
                       L_0=2\cdot80989^2=13118436242.
\]
The complete integer comparison verifies
\[
 80989^{94}+1>C_0,\qquad80989^{47}>L_0,
 \qquad80989^{90}+1<C_0.
\]
Consequently, for both families,
\[
                        H(P_N)=80989^{4N+2}+1
                        \qquad(N\ge23).
\]
Indeed this coefficient exceeds both $C_0$ and $L_0A_N$ at $N=23$,
and continues to do so as $N$ increases. Since $C_0>L_0^2$, the exact
height for any positive parameter is $\max(C_0,A_N^2+1)$.
This is an algebraic height threshold; the sufficient parameter in the
342-variable defining theorem remains $27464$, and the corresponding
threshold for the 297-variable family remains ineffective.

\section{Soundness for every sufficiently large parameter}
Put $\beta=51X^2+2504$ and $c=|a-1|$. The retained scalar circuit
uses the positive contact product
\[
 r=\beta(\beta-1)(\beta-2401)(32\beta+343)(2\beta-1)
\]
and normalizer $R_s=A_Nr/(1+80989^2r)$, with shift $R_s+9792$.
The fixed-prime argument gives
\begin{equation}\label{eq:scalar}
 c\ge80989^{L_N},\quad
 L_N=\left\lfloor\frac{187(2N-1)}8\right\rfloor+1,\quad
 h(\beta)\le\HH c^{19/24}.
\end{equation}
Its integral group identities and exceptional-point exclusions are
unaffected by the local test factors.

\begin{proof}[Proof of soundness in Theorem~\ref{thm:main}]
Set
\[
 \Delta=t^2-a(a-2)tb+(a-2)^2b^2-1.
\]
The formal comparison gives divisibility of $\Delta$ by the fortieth
powers of both contact radicals. If $\Delta\ne0$, then
$|\Delta|\le14B^4$, and \eqref{eq:height} gives
\[
                     B^2\le C14^{19/40}B^{19/10}+M_2.
\]
This is impossible for sufficiently large $B$. Equation~\eqref{eq:scalar}
ensures $B\ge c-1\to\infty$ uniformly as $N\to\infty$. Hence
$\Delta=0$ for every sufficiently large $N$.

The integral Lucas descent in \cite{Source} then gives
$\log B>|a|-3\ge c-4$. Suppose $n-\beta\ne0$, and let $J$ be
its nonzero reduced numerator. The order-ten tests imply that both
tenth powers of the contact radicals divide $J$, while
$|J|\le2B\exp(h(\beta))$. For $B^2\ge2M_2$, equation~\eqref{eq:height}
therefore implies
\[
 \log B\le10\log(2C)+19\log2+19h(\beta)
          \le10\log(2C)+19\log2+19\HH c^{19/24}.
\]
The right side is $o(c)$, contradicting the Lucas lower bound for
sufficiently large $c$. Thus $n=\beta\in\Z$. Writing $X=u/v$ in
lowest terms, $v>0$, gives $v^2\mid51$ and hence $v=1$.

All thresholds depend only on the fixed constants $C,M_2,\HH$.
Since $L_N$ increases, one integer $N_0\ge27464$ works for every
$N\ge N_0$. The dependence on $M_2$ is precisely the ineffective
part of the assertion.
\end{proof}

\section{Completeness and exact verification}
Fix an integer input and the scalar parameter. The retained scalar and
Lucas construction supplies the integer labels and nonzero root points.
Choose compact smaller real and finite-place neighborhoods avoiding the
poles of the retained rational tests. The tests then have fixed valuation
lower bounds at every prime in the finite preparation set $S$.
All scales and neighborhoods are fixed before the contact primes.

Choose a positive integer $D$ that clears the contact coefficients
$D^2\theta/a_*$ and has sufficiently large valuations at the primes in
$S$. In the simultaneous prime pattern put
\[
 \lambda=\frac{u}{D^2w^2}.
\]
The zero contact makes $u$ a selected prime, outside $S$. Therefore
$v_p(\lambda)\le-2v_p(D)$. Making this upper bound sufficiently negative
gives $v_p(t_*)<\min(0,v_p(\theta):\theta\ne0)$ and hence
\[
 v_p(F_*(t_*))=5v_p(t_*),\quad
 v_p(e)=5v_p(a_*),\quad
 v_p(q)=2v_p(a_*)-3v_p(\lambda).
\]
The latter is uniformly arbitrarily large. Thus all new weighted test
terms are arbitrarily small at the prepared finite places. Enlarging $D$
is a fixed dilation of the prime pattern and is allowed by the retained
selection theorem; see \cite{Selection}.

At each selected contact prime, the designated primary place gives
all test orders. The distinct-valuation argument at the opposite place
makes $\Tr(W^{-1})$ a unit congruent to $1$ or $2$ modulo $p>2$.
Consequently the pole input $R_*\Tr(W^{-1})^{-2}$ has valuation one.
The exact product formula
\[
 R_*=\frac{\prod_{\theta\in\Theta}(u-(\theta/a_*)D^2w^2)}
               {(D^2w^2)^5\delta^6}
\]
shows that its valuation is even at every other prime. This includes
all primes dividing $\delta$; no unit assertion about $\delta$ away from
contacts is required. Its reduced denominator is therefore a square.
The selected-prime congruences give the required local squareclasses
at $2,3,17$, and the fixed choice of $D$ gives their required valuations.

Uniformly on the fixed compact real neighborhoods,
\[
 R_*\sim a_*^{-6}\lambda^{-1},\qquad q\sim a_*^2\lambda^{-3},
 \qquad \Tr(W^{-1})\longrightarrow2.
\]
All factors tend to $1$, and the pole input tends to $0^+$. It therefore
lies in $(0,17)$ for sufficiently large $\lambda$. The simultaneous
selection supplies such parameters after all fixed data have been chosen.
The pole-realization premise gives fresh witnesses, and the rational
weight identities fill every new quadratic row. Equal scales may reuse
witnesses. This is the full simultaneous completeness argument of
\cite[Proposition 3.2]{Shared}; it imposes no new lower bound on the
scalar parameter $N$.

The exact reproduction programs are
\begin{verbatim}
python build_quartic_shared_weights.py
python check_quartic_shared_weights.py
\end{verbatim}
The checker imports no producer. It checks every changed and retained
row, all coordinate incidences, and every coefficient of the quartic by
an independent unordered-pair expansion. It verifies degree four,
content one and the coefficient-height bound. These finite checks do
not prove the retained global arithmetic inputs.

The exact table contains 271 rows, 297 total coordinates, 4269 terms
before parameter powers are collected, and 4268 ordinary monomials for
every $N\ge1$. The weight presentation uses one coordinate $c$ and one coordinate
$q_8$ in each detector, together with the five displayed power outputs.
The coefficient family is fully specified; the sufficient defining
threshold remains ineffective because of the aggregate trace height.

\begingroup
\small
\begin{thebibliography}{9}
\bibitem{Shared}
\emph{Common local weights and quartics in 297 and 342 variables},
9 November 2026, Sections 2--4: local extraction, simultaneous preparation,
nonzero domains and exact quadratic weight presentation.
\href{../local_tests/shared_local_weights_20261109_v12.pdf}{PDF};
source: \nolinkurl{shared_local_weights_20261109_v12.tex}.
\bibitem{Selection}
\emph{Separation and denominator conditions in quadratic prime patterns},
9 November 2026, weighted separation from Sun's stated analytic estimates.
\href{../prime_selection/prime_selection_20261109_v12.pdf}{PDF};
source: \nolinkurl{prime_selection_20261109_v12.tex}.

\bibitem{Packing}
\emph{Simultaneous valuation conditions and explicit quartics},
9 November 2026, accompanying local theorem, simultaneous completeness
proof, exact circuit, and independent verification.
\href{../local_tests/mixed_valuation_tests_20261109_v11.pdf}{PDF}; source: \nolinkurl{mixed_valuation_tests_20261109_v11.tex}.
\bibitem{Explicit}
\emph{An explicit quartic in 342 variables}, 9 November 2026,
specified polynomial and effective conditional projection theorem.
\href{../quartic/quartic342_20261109_v12.pdf}{PDF}; source: \nolinkurl{quartic342_20261109_v12.tex}.
\bibitem{Baseline}
\emph{A family of quartics in 317 variables from aggregate trace heights},
8 November 2026, full point-chart, trace-height, and scalar arguments.
\href{quartic317_family_20261108_v9.pdf}{PDF}; source: \nolinkurl{quartic317_family_20261108_v9.tex}.
\bibitem{Source}
\emph{A forced prime condition and an explicit quartic in 392 variables},
8 November 2026, arithmetic input contract and scalar/Lucas arguments.
\href{../quartic/quartic392_20261108_v8.pdf}{PDF}; source: \nolinkurl{quartic392_20261108_v8.tex}.
\bibitem{Aggregate}
\emph{Aggregate trace heights of elliptic multiples},
8 November 2026, proof of the linear height comparison from Von Buhren's
family height theorem.
\href{../trace_heights/aggregate_trace_height_20261108_v9.pdf}{PDF}; source: \nolinkurl{aggregate_trace_height_20261108_v9.tex}.
\bibitem{Glazer}
E.~Glazer, \emph{An explicit quartic defining the integers in the rationals},
Version~13, manuscript and coefficient data, commit
\nolinkurl{5bbd74f6a479776a332cbe036809d5dddae63720},
\url{https://github.com/elliotglazer/diophantine-z-q}.
\end{thebibliography}
\endgroup
\end{document}
